\fancypagestyle{plain}{
    \fancyhf{} % 清除所有页眉页脚重新设置
    \fancyhead[C]{\songti\zihao{5}天津大学20XX届本科生毕业论文} % 页眉，宋体五号
    \fancyfoot[C]{\songti\zihao{-5} \thepage} % 页脚，宋体小五号，\thepage 代指页码
    \renewcommand{\headrulewidth}{0.5bp} % 页眉横线粗细
    \renewcommand{\footrulewidth}{0bp} % 页脚横线粗细
}

\fancypagestyle{noheader}{
    \fancyhf{} % 清除所有页眉页脚重新设置
    \fancyfoot[C]{\songti\zihao{-5} \thepage} % 页脚，宋体小五号，\thepage 代指页码
    \renewcommand{\headrulewidth}{0bp} % 页眉横线粗细
    \renewcommand{\footrulewidth}{0bp} % 页脚横线粗细
}

\makeatletter

\renewcommand\frontmatter{
    \@mainmatterfalse
    \clearpage
    \pagestyle{noheader} % 使用页眉页脚样式 noheader
    \pagenumbering{Roman} % 大写罗马数字页码（会重置页码）
    \setlength{\baselineskip}{20bp plus 2bp minus 2bp} % 设置弹性行高为 20±2bp （optional）
} 
\renewcommand\mainmatter{
    \@mainmattertrue
    \clearpage
    \pagestyle{plain} % 使用页眉页脚样式 plain 
    \pagenumbering{arabic} % 使用阿拉伯数字页码（会重置页码）
} 

\makeatother

\ctexset{
    contentsname={目\quad 录}, % 目录的标题 
    secnumdepth=3, % 标题编号深度
    chapter={
        format={\centering\zihao{-3}\heiti},   % 小三，黑体，居中
        number={\chinese{chapter}}, % 章节编号黑体中文数字
        name={第,章},   % 章前缀
        aftername=\quad, 
        beforeskip=30bp, % 段前间距 （规定 30~36 bp 之间）
        afterskip=30bp, % 段后间距 （规定 30~36 bp 之间）
    },
    section={
        format={\zihao{4}\heiti}, % 四号黑体
        beforeskip=20bp, % 段前间距 （规定 18~24 bp 之间）
        afterskip=20bp, % 段后间距 （规定 18~24 bp 之间）
        aftername=\quad, % 编号与标题之间空格
    },
    subsection={
        format={\zihao{4}\heiti}, % 四号黑体
        beforeskip=12bp, % 段前间距 （规定 12~15 bp 之间）
        afterskip=12bp, % 段后间距 （规定 12~15 bp 之间）
        aftername=\quad,
    }
}

\addtocontents{toc}{\protect\thispagestyle{noheader}} % 设置目录页的页眉页脚样式为 noheader 
\setcounter{tocdepth}{2} % 设置 tocdepth 目录深度 （2 = subsection，即三级标题）
\titlecontents{chapter}
    [0em] % 缩进为 0em = 0 字符
    {\heiti\zihao{-4}}
    {\thecontentslabel\quad}{}
    {\hspace{.5em}\titlerule*[5bp]{.}\zihao{-4}\contentspage} %
\titlecontents{section}
    [2em] % 缩进为 2em = 4 字符
    {\zihao{-4}}
    {\thecontentslabel\quad}{}
    {\hspace{.5em}\titlerule*[5bp]{.}\zihao{-4}\contentspage}
\titlecontents{subsection}
    [4em] % 缩进为 4em = 4 字符
    {\zihao{-4}}
    {\thecontentslabel\quad}{}
    {\hspace{.5em}\titlerule*[5bp]{.}\zihao{-4}\contentspage}

\renewcommand{\tablename}{表}                                     % 插表题头
\renewcommand{\figurename}{图}                                    % 插图题头
\renewcommand{\thefigure}{\arabic{chapter}-\arabic{figure}}       % 使图编号为 7-1 的格式 %\protect{~}
% \renewcommand{\thesubfigure}{\alph{subfigure})}                   % 使子图编号为 a) 的格式
% \renewcommand{\thesubtable}{(\alph{subtable})}                    % 使子表编号为 (a) 的格式
\renewcommand{\thetable}{\arabic{chapter}-\arabic{table}}         % 使表编号为 7-1 的格式
\renewcommand{\theequation}{\arabic{chapter}-\arabic{equation}}   % 使公式编号为 7-1 的格式
% \newcommand{\ud}{\mathrm{d}} % 定义微分符号的命令，使用 upright 字体

\newcommand{\figuremacro}[4]{
  \begin{figure}[#1]
    \centering
    \includegraphics[width=#4\textwidth]{#2}
    \caption{\zihao{5}#3}
    \label{fig:\thefigure}
  \end{figure}
}
